\scanpage{92}
\deleted{Si \unsure{dans}}

Si dans les catégories fibres, certains \inserted{types de} $\underrightarrow{\lim}$ ou $\underleftarrow{\lim}$
sont représentables, et si les $f^*$ y commutent,
alors les $=\mathrm{s}$ types de lim sont représentables
dans $\underleftarrow{\mathrm{Lim}}\ \unsure{E|\mathcal E}/B$, et les foncteurs fibres (plus
gén, les foncteurs chang\textsuperscript{ts} de base par \unsure{des|un}
$B'\mathbin{\unsure{-|\to}}B$) y commutent. Ainsi, si les $E_b$\unsure{/|{}}
sont préadditives (\unsure{r.|resp.} additives, abéliennes) \inserted{et si les $f^*$ ``respectent'' ces struct.} il en est
de $=$ de \deleted{$\unsure{\Gamma}$} $\underleftarrow{\mathrm{Lim}}\ E/B$, si $E_b$ est à \inserted{petites} $\underrightarrow{\lim}$ filtrantes
\unsure{ex.ctes|exactes}, et si les $\underleftarrow{\lim}$ finies y sont représentables,
\inserted{\uline{et} si les $f^*$ commutent à \unsure{tt ça|t.t. ça}, alors} kif-kif pour $\underleftarrow{\mathrm{Lim}}\ E/B$. De $=$, si les $E_b$
satisfont \inserted{\unsure{l'une}} les conditions a) b) c) \unsure{de l'exactitude|d'exactitude} d'
un $\mathcal U$-topos, \inserted{et les $f^*$ des morphismes de topos,} de $=$ $\underleftarrow{\mathrm{Lim}}\ E/B$. Pour pouvoir
conclure que si les $E_b$ sont des $\mathcal U$-topos
(resp. des $\mathcal U$-topos \unsure{abéliens|distribués}) \inserted{et les $f^*$ des morphismes de tels,} alors il en est de
$=$ de \deleted{$\unsure{\Gamma}$} $\underleftarrow{\mathrm{Lim}}\ E/B$, il reste à donner des
conditions assurant que cette catégorie est une
$\mathcal U$-catégorie, et admet une \inserted{petite} famille de gén.
Le premier pt sera assuré dès que $B$ est
``ess. petite''. Pour \unsure{le|les} deuxième, si on veut
appliquer le corollaire qui précède, on est

% archive sheet number: 91
